%! AP Statistics — Unit 3, Lesson 3.3 — Student Guided Notes
\documentclass[10pt]{article}
\usepackage{apstats-article}
\usepackage{apstats-boxes}

\begin{document}

\pageheader{Unit 3, Lesson 3.3}{Guided Notes}
\namedateperiod

\vspace{0.12in}

\begin{objectivebox}
\small By the end of this lesson, I will be able to\dots\\[4pt]
\writeline\\[1pt]
\writeline\\[1pt]
\writeline
\end{objectivebox}

\vspace{0.1in}

\begin{vocabbox}
\termblanklong{Simple random sample (SRS)}
\termblanklong{Stratified random sample}
\termblanklong{Cluster sample}
\termblanklong{Systematic random sample}
\termblanklong{Census}
\end{vocabbox}

\vspace{0.1in}

\begin{notesbox}{Sampling Methods Reference}
\small
\renewcommand{\arraystretch}{2.0}
\begin{tabularx}{\linewidth}{@{} >{\bfseries\color{navy}}p{0.24\linewidth}
    >{\raggedright\arraybackslash}p{0.38\linewidth}
    >{\raggedright\arraybackslash}X @{}}
\textbf{Method} & \textbf{How to do it} & \textbf{Best when \ldots} \\
\hline
Simple Random Sample &
Number all individuals; randomly select $n$ using a random device. Every group of
size $n$ has \blank{2cm} chance of being chosen. &
The population is \blank{3cm} \\
Stratified Random Sample &
Divide into \blank{2cm} groups called \blank{2cm} (homogeneous within).
SRS within each stratum. &
You want to ensure representation of \blank{3cm} \\
Cluster Sample &
Divide into \blank{2cm} (heterogeneous). Randomly select clusters; survey
\blank{3cm} in selected clusters. &
Individuals are \blank{3cm} geographically \\
Systematic Random Sample &
Choose a \blank{3cm} start; select every \blank{2cm}th individual after. &
You have a complete \blank{3cm} list \\
Census &
Measure \blank{5cm} in the population. &
Population is \blank{3cm} and manageable \\
\end{tabularx}
\end{notesbox}

\vspace{0.1in}

\begin{notesbox}{Stratified vs. Cluster --- KEY Distinction}
\small
\begin{tabularx}{\linewidth}{@{} >{\bfseries\color{navy}}p{0.3\linewidth} X @{}}
Strata are \ldots & \blank{5cm} (same type within each group) \\
Clusters are \ldots & \blank{5cm} (mixed types within each group) \\
In stratified: survey & \blank{3cm} from \blank{3cm} stratum \\
In cluster: survey & \blank{3cm} from \blank{3cm} selected clusters \\
\end{tabularx}

\vspace{8pt}
\textbf{Memory trick:} Strata = \textbf{S}ame within. Clusters = mi\textbf{X}ed within.
\end{notesbox}

\vspace{0.1in}

\begin{practicebox}
\small Identify the sampling method in each scenario.
\begin{enumerate}[leftmargin=1.5em, itemsep=10pt]
  \item A researcher numbers all 2{,}000 employees at a company and uses a random
        number table to select 100. \blank{7cm}

  \item A school divides students into 4 grade levels, then randomly selects 25
        students from each grade. \blank{7cm}

  \item A city divides its 60 neighborhoods into groups; randomly selects 6
        neighborhoods; surveys every household in the selected neighborhoods.
        \blank{7cm}

  \item A manager numbers all 500 customers 1--500, randomly picks a start between
        1--10, then surveys every 10th customer. \blank{7cm}
\end{enumerate}
\end{practicebox}

\end{document}
